\scanpage{41}
on a alors
\[
T_{\dot g}i=(\mathrm{det}(\dot g))^\bullet.i\tag*{(98)}
\]
ce qui relie (93) et (97), et caractérise en fait (97) en termes de (93). Pour $\dot g$ de la forme $\dot\lambda$, avec $\dot\lambda\in Z'_{\unsure{\Lambda|M}}$, on trouve donc
\[
T_{\dot\lambda}i=\dot\lambda^2.i\qquad\dot\lambda\in Z'_\Lambda,\quad i\in\pi_0(P_M),\tag*{(99)}
\]
formule qui relie les deux éléments de structure (95) et (97).

Enfin, rappelons nous des bijections canoniques
\[
\deleted{\unsure{P}}S_{M,i}\simeq M^*\qquad i\in\pi_0(P_M),\tag*{(100)}
\]
où $\deleted{\unsure{P}}S_{M,i}$ désigne la \inserted{l'ens des sommets de la} composante connexe $P_{M,i}$ de $P_M$ d'indice $i$ -- ces bijections canoniques définissent \deleted{\unsure{une}},
\[
\left\{\begin{aligned}
&\text{un syst. transitif de bijections canoniques}\\
&\mathrm{tr}_{\unsure{s|S}\,ji}:S_{M,i}\overset{\sim}{\text{---}} S_{M,j}
\qquad(i,j\in\pi_0(P_M)).
\end{aligned}\right.\tag*{(101)}
\]
Appelons \uline{automorphisme strict} de $P_M$ tout automorphisme \deleted{\unsure{qui conserve}}\inserted{du polyèdre qui} respecte les trois structures supplémentaires (95) (97) (101).
\margin{On verra par la suite que les \unsure{deux} structures (95) (97) sont inutiles}
% slanted
Désignons \deleted{alors \unsure{avec}} \inserted{par $\mathrm{Autstr}$}$(P_M)$
% archive sheet number: 40
